\section{The strong barrier}
\label{sec:barrier}

This section is elementary and contains no small parameter. The inequality $C\ge\sum_i\Lambda_i\Dfn(e_i)$ turns the strong condition $C\le B(\Lambda)-m$ into a bound on the eccentricities that keeps every perihelion and every pair of adjacent radial annuli uniformly apart, and $B$ is the least deficit at which two annuli can touch. Conservation of $L$ then makes $C$ move only with the actions (Lemma~\ref{lem:moving}). Throughout, $d\in\{2,3\}$, and in the plane $q\times p$ is the oriented scalar determinant.

\begin{lemma}[Normalization]
\label{lem:normalization}
For every $d\in\{2,3\}$ and all $\nu_i,\mu>0$, the Hamiltonian \eqref{eq:H} is the translation reduction, at zero total momentum, of the Newtonian $(1+n)$-body Hamiltonian with masses $(1,\mu\nu_1,\dots,\mu\nu_n)$, divided by $\mu$, with unchanged time. Its solutions reconstruct Newtonian solutions with these masses.
\end{lemma}

\begin{proof}
In barycentric coordinates put $P_i=\mu p_i$, $P_0=-\sum_iP_i$ and $q_i=x_i-x_0$, so that $\sum_{i=0}^nP_i\cdot dx_i=\sum_{i=1}^nP_i\cdot dq_i$. The kinetic energy is $\sum_i\abs{P_i}^2/(2\mu\nu_i)+\frac12\abs{\sum_iP_i}^2$ and the potential $-\sum_i\mu\nu_i/\abs{q_i}-\sum_{i<j}\mu^2\nu_i\nu_j/\abs{q_i-q_j}$; the symplectic form and the Hamiltonian both acquire the factor $\mu$, which cancels and gives \eqref{eq:H}. Conversely, put $x_0=-\mu\sum_i\nu_iq_i/(1+\mu\sum_i\nu_i)$ and $x_i=x_0+q_i$; Hamilton's equations give $\dot x_0=-\mu\sum_jp_j$, $\dot x_i=p_i/\nu_i$, and Newton's equations for the momenta $\mu p_i$.
\end{proof}

\begin{lemma}[Orbital identities]
\label{lem:orbital}
Let $k>0$, $q\ne0$, $v\in\R^d$ with $E=\abs{v}^2/2-k/\abs{q}<0$, and put $a=-k/(2E)$, $\ev=\big((\abs v^2-k/\abs q)q-(q\cdot v)v\big)/k$ and $e=\abs\ev$. Then, including the radial case $e=1$,
\begin{equation}
\label{eq:orb}
e^2=1+\frac{2E\abs{q\times v}^2}{k^2}=1-\frac{\abs{q\times v}^2}{ka}\le1,\qquad
a(1-e)\le\abs{q}\le a(1+e),\qquad \abs{v}^2=k\Big(\frac{2}{\abs{q}}-\frac1a\Big).
\end{equation}
For $k=1$ and $v=p/\nu$ these are the elements \eqref{eq:elements}, and $\abs{q\times p}=\Lambda\sqrt{1-e^2}$.
\end{lemma}

\begin{proof}
Expanding $k^2e^2$ and using $\abs{q\times v}^2=\abs q^2\abs v^2-(q\cdot v)^2$ gives the first identity; substitute $2E=-k/a$ for the second. With $r=\abs q$, the energy identity gives $(q\cdot v)^2=2kr-kr^2/a-\abs{q\times v}^2\ge0$; inserting $\abs{q\times v}^2=ka(1-e^2)$ and completing the square gives $(r-a)^2\le a^2e^2$. The last identity is the definition of $E$, and $\abs{q\times p}=\nu\abs{q\times v}=\nu\sqrt{a(1-e^2)}$ for $k=1$.
\end{proof}

\begin{lemma}[Angular momentum and the AMD identity]
\label{lem:amd}
Along every solution of \eqref{eq:H} on an interval, $L$ is constant. In the negative-energy region $C\ge\sum_i\Lambda_i\Dfn(e_i)\ge0$. If $C<\Bst(\Lambda)$, then $L\ne0$, the projections $H_i=L_i\cdot L/\abs{L}$ are positive, and
\begin{equation}
\label{eq:amd}
C(t)-C(0)=\sum_i\big(\Lambda_i(t)-\Lambda_i(0)\big),\qquad C=\sum_i(\Lambda_i-H_i).
\end{equation}
If $C\le\Bst(\Lambda)-m$ with $m>0$, then $H_i\ge m$ for every $i$. In the plane, after choosing the orientation with $L>0$, every $L_i$ is positive and $H_i=L_i$.
\end{lemma}

\begin{proof}
In $\frac{d}{dt}\sum_iq_i\times p_i$ the kinetic terms $(p_i/\nu_i)\times p_i$ and $\mu(\sum_jp_j)\times(\sum_ip_i)$ vanish, central forces have zero torque, and the two forces of a pair give a multiple of $(q_i-q_j)\times(q_i-q_j)=0$; so $\dot L=0$. By \eqref{eq:orb}, $\abs L\le\sum_i\abs{L_i}=\sum_i\Lambda_i\sqrt{1-e_i^2}$, which gives $C\ge\sum_i\Lambda_i\Dfn(e_i)\ge0$. If $L=0$, then $C=\sum_i\Lambda_i\ge\Bst$. Otherwise $\abs L=\sum_iH_i$ gives the second identity in \eqref{eq:amd}, and constancy of $\abs L$ gives the first. Since each $\Lambda_i-H_i\ge\Lambda_i-\abs{L_i}\ge0$, we get $H_i\ge\Lambda_i-C$, which is positive if $C<\Bst$ and at least $m$ if $C\le\Bst-m$.
\end{proof}

\begin{lemma}[The thresholds]
\label{lem:threshold}
For all $\nu_i>0$ and $\Lambda\in\Ach$, every minimum in \eqref{eq:thresholds} is attained and positive, and $B_{ij}$ and $B$ are continuous on $\Ach$. The minimum defining $B_{ij}$ is attained at radial contact $a_i(1+u)=a_j(1-v)$, and
\begin{equation}
\label{eq:Bsharp}
B_{ij}(\Lambda)<\Lambda_j,\qquad B(\Lambda)=\min\big(\Bcross(\Lambda),\Lambda_1\big).
\end{equation}
With $\alpha=a_i/a_j$, the quotient $B_{ij}(\Lambda)/\Lambda_j$ is the relative critical AMD $C_c(\alpha,\nu_i/\nu_j)$ of Laskar and Petit~\cite[\S\S3.1--3.2]{LP17}.
\end{lemma}

\begin{proof}
Fix $i<j$, so $0<\alpha<1$. The feasible set $\{(u,v)\in[0,1]^2:\alpha u+v\ge1-\alpha\}$ is compact and contains $(0,1-\alpha)$, and the objective $\Lambda_i\Dfn(u)+\Lambda_j\Dfn(v)$ is continuous, so the minimum is attained. If $(u,v)$ is feasible, then $t=(1-\alpha)/(\alpha u+v)\le1$, the point $(tu,tv)$ lies on the contact line $\alpha u+v=1-\alpha$, and since $\Dfn$ is increasing the cost does not increase. Hence the minimum is attained at contact, which we parametrize over a fixed interval:
\begin{equation}
\label{eq:contact}
u=\bar u_\alpha t,\qquad v=1-\alpha-\alpha\bar u_\alpha t,\qquad t\in[0,1],\qquad \bar u_\alpha=\min\Big(1,\frac{1-\alpha}{\alpha}\Big).
\end{equation}
The resulting cost $F(\Lambda,t)$ is continuous, and $\abs{\min_tF(\Lambda,t)-\min_tF(\Lambda',t)}\le\sup_t\abs{F(\Lambda,t)-F(\Lambda',t)}$ with uniform continuity on compact sets proves continuity of $B_{ij}$, hence of $B$. A zero cost forces $u=v=0$, which is not feasible. The competitor $(0,1-\alpha)$ gives $B_{ij}\le\Lambda_j\Dfn(1-\alpha)<\Lambda_j$. In particular $\Lambda_i>B_{i-1,i}\ge\Bcross$ for $i\ge2$, which gives the second identity in \eqref{eq:Bsharp}.

Finally $\Lambda_i/\Lambda_j=(\nu_i/\nu_j)\sqrt\alpha$, so $B_{ij}/\Lambda_j$ is the minimum of $(\nu_i/\nu_j)\sqrt\alpha\,\Dfn(u)+\Dfn(v)$ on the contact line $\alpha u+v=1-\alpha$, $u,v\in[0,1]$. This is the minimization of the relative AMD under the collision condition that defines $C_c$ in~\cite[\S\S3.1--3.2]{LP17}, with mass ratio $\nu_i/\nu_j$.
\end{proof}

The threshold is geometrically sharp. If $(u,v)$ minimizes $B_{ij}$, two coplanar prograde ellipses with eccentricities $u,v$ and opposite perihelion directions touch, the inner aphelion meeting the outer perihelion, and their combined deficit is $B_{ij}$; a radial orbit has deficit $\Lambda_i$ and is the limit of prograde ellipses as $e_i\uparrow1$. This is not a claim that motion above the threshold must collide.

\begin{lemma}[Uniform barrier]
\label{lem:uniform}
Let $U\subset\Ach$ be compact and $m>0$. There exist $\dstar,\dpair>0$, depending only on $(n,\nu,U,m)$, such that for every $\Lambda\in U$ and $e\in[0,1]^n$ with $\sum_i\Lambda_i\Dfn(e_i)\le B(\Lambda)-m$,
\begin{equation}
\label{eq:barrier}
a_i(\Lambda)(1-e_i)\ge\dstar,\qquad a_{i+1}(\Lambda)(1-e_{i+1})-a_i(\Lambda)(1+e_i)\ge\dpair.
\end{equation}
Every state with $q_i\ne0$, $E_i<0$, $\Lambda\in U$ and $C\le B(\Lambda)-m$ satisfies \eqref{eq:barrier}, $\abs{q_i}\ge\dstar$ and $\abs{q_i-q_j}\ge\dpair$ for $i\ne j$, and $H_i\ge m$; these states form a compact subset of the region $\{q_i\ne0,\ q_i\ne q_j,\ E_i<0\}$.
\end{lemma}

\begin{proof}
Let $K_m=\{(\Lambda,e)\in U\times[0,1]^n:\sum_i\Lambda_i\Dfn(e_i)\le B(\Lambda)-m\}$, which is compact by Lemma~\ref{lem:threshold}. If an adjacent separation in \eqref{eq:barrier} were nonpositive at a point of $K_m$, then $(e_i,e_{i+1})$ would be feasible for $B_{i,i+1}$ and $\sum_k\Lambda_k\Dfn(e_k)\ge B_{i,i+1}(\Lambda)\ge B(\Lambda)$, a contradiction. So the minimum $\dpair$ of the continuous separations over $K_m$ is positive (if $K_m=\emptyset$ all is vacuous). Put
\begin{equation}
\label{eq:dstar}
\Lambda_{\max}=\max_{\Lambda\in U,\,i}\Lambda_i,\qquad a_{\min}=\min_{\Lambda\in U,\,i}a_i(\Lambda),\qquad a_{\max}=\max_{\Lambda\in U,\,i}a_i(\Lambda),\qquad \dstar=\frac{a_{\min}m^2}{2\Lambda_{\max}^2}.
\end{equation}
As $B\le\Bst\le\Lambda_i$, the deficit bound gives $\Lambda_i\Dfn(e_i)\le\Lambda_i-m$, i.e.\ $\sqrt{1-e_i^2}\ge m/\Lambda_{\max}$; with $1-e_i\ge(1-e_i^2)/2$ this gives the first bound in \eqref{eq:barrier}.

For a physical state, Lemma~\ref{lem:amd} gives $\sum_i\Lambda_i\Dfn(e_i)\le C\le B(\Lambda)-m$, so \eqref{eq:barrier} holds, and $H_i\ge m$ because $B\le\Bst$. By \eqref{eq:orb}, $\dstar\le\abs{q_i}\le2a_{\max}$, $\abs{p_i}\le\nu_i\sqrt{2/\dstar}$ and $E_i\le-1/(2a_{\max})$, and summing adjacent separations gives $\abs{q_i-q_j}\ge\abs{q_j}-\abs{q_i}\ge\dpair$ for $i<j$. These inequalities define a compact set inside the nonsingular region, on which $\Lambda\in U$ and $C\le B(\Lambda)-m$ are closed conditions.
\end{proof}

Theorem~\ref{thm:barrier} is Lemma~\ref{lem:uniform}: in the plane, with the orientation $L>0$ of Lemma~\ref{lem:amd}, $H_i=L_i=\abs{L_i}\ge m$. The next lemma, the only place where conservation of $L$ enters, converts action drift into margin drift.

\begin{lemma}[Moving threshold and continuation]
\label{lem:moving}
Let $\Lset\subset\Ach$ be compact and $\delta>0$. There exist a radius $r_g>0$ and a compact $U\subset\Ach$ containing the closed $\nrm{\cdot}$-neighbourhood of radius $4r_g$ of $\Lset$ such that the modulus
\[
\Omega_B(r)=\sup\{\abs{B(\Lambda)-B(\Lambda')}:\Lambda,\Lambda'\in U,\ \nrm{\Lambda-\Lambda'}\le r\}
\]
satisfies $\Omega_B(r)\to0$ as $r\downarrow0$ and $\Omega_B(4r_g)\le\delta/16$. Let $0<r\le4r_g$. If a solution of \eqref{eq:H} has $\Lambda^0=\Lambda(0)\in\Lset$, $C(0)\le B(\Lambda^0)-\delta$, and $\nrm{\Lambda(t)-\Lambda^0}\le r$ while $E_i(t)<0$, then
\begin{equation}
\label{eq:moving}
C(t)\le B(\Lambda(t))-\delta+nr+\Omega_B(r).
\end{equation}
If $nr+\Omega_B(r)\le\delta/2$, the maximal solution cannot end, forwards or backwards, at a finite time while these conditions hold.
\end{lemma}

\begin{proof}
Take $\rho>0$ with the closed $2\rho$-neighbourhood $U$ of $\Lset$ inside $\Ach$. By Lemma~\ref{lem:threshold}, $B$ is uniformly continuous on $U$, so $\Omega_B(r)\to0$; choose $4r_g\le\rho$ with $\Omega_B(4r_g)\le\delta/16$. Lemma~\ref{lem:amd} gives $C(t)=C(0)+\sum_i(\Lambda_i(t)-\Lambda_i^0)\le B(\Lambda^0)-\delta+nr$, and $B(\Lambda^0)\le B(\Lambda(t))+\Omega_B(r)$ gives \eqref{eq:moving}. If $nr+\Omega_B(r)\le\delta/2$, then $C(t)\le B(\Lambda(t))-\delta/2$, and Lemma~\ref{lem:uniform} with $m=\delta/2$ keeps the state in a fixed compact subset of the domain of the analytic vector field of \eqref{eq:H}, which a maximal solution must leave at a finite endpoint, in either time direction.
\end{proof}
